\chapter{Introducción} \label{ch:introduccion}

% --- PÁRRAFO 1: FUNDAMENTOS DINÁMICOS Y COMPLEJIDAD DEL FENÓMENO ---

Desde una perspectiva dinámica global, el análisis de los storm tracks realizado por \citet{hoskins2005new} nos muestra a los ciclones extratropicales como el mecanismo dominante para el transporte latitudinal de energía en el Hemisferio Sur. La cuantificación de estos transportes ha sido posible gracias al desarrollo de algoritmos de seguimiento objetivo, donde \citet{sinclair1994objective} estableció el precedente fundamental al utilizar la vorticidad relativa para identificar y rastrear sistemas independientemente del flujo de fondo, permitiendo entender que estos sistemas equilibran el gradiente térmico planetario.

En el contexto regional, \citet{reboita2010south} proporcionaron una de las primeras climatologías sistemáticas de alta resolución para el Atlántico Sur, identificando tres centros preferenciales de ciclogénesis vinculados a la interacción entre el flujo perturbado por los Andes y los gradientes térmicos oceánicos. 

En el litoral sur de Brasil, este incremento se traduce en riesgos meteoceanográficos severos; como documentan recientemente \citet{miranda2026patterns}, los ciclones extratropicales son los principales motores de tormentas costeras en esta región.

No obstante, la naturaleza de estas perturbaciones en la región es heterogénea. \citet{evans2012climatology} destacan la prevalencia de ciclones subtropicales con estructuras térmicas híbridas. Esta complejidad desafía la clasificación binaria tradicional, lo que concuerda con \citet{hart2003cyclone}, quien demostró que los sistemas transitan a lo largo de un continuo estructural en lugar de pertenecer a categorías estáticas. Complementariamente, \citet{gozzo2013air} documentan la ocurrencia de sistemas que siguen el modelo estructural de Shapiro-Keyser, caracterizados por una marcada fractura del frente frío y el desarrollo de una seclusión cálida durante su etapa de madurez. Estas variaciones morfológicas, respaldadas por climatologías regionales de largo plazo \citep{gozzo2014subtropical}, evidencian que el Atlántico Sur se aparta frecuentemente de la conceptualización de ciclones convencionales, lo que subraya la idea de caracterizar cómo la arquitectura interna de cada sistema modula la distribución espacial de las estructuras de viento y precipitacion en superficie. Así, más allá de su rol climático, en el hemisferio sur, estos sistemas representan el fenómeno de severidad geofísica dominante para la zona costera regional \citep{gramcianinov2023impact}.

Esta diversidad estructural tiene implicancias directas en cómo medimos su severidad: como argumenta \citet{corner2025classification} para el Atlántico Norte, la intensidad de un ciclón no puede capturarse con una única métrica, por que la relación entre la intensidad del sistema y los vientos en superficie no es lineal. De hecho, \citet{eisenstein2023identification} destacan que los mayores impactos en superficie (en ciclones del hemisferio norte) suelen estar asociados a estructuras de mesoescala específicas, como los cinturones de transporte frío (CCB, por su siglas en ingles), cuya localización espacial respecto al centro de vorticidad varía drásticamente durante el ciclo de vida del sistema. 

% --- PÁRRAFO 2: EL "GAP" ESTRUCTURAL Y LA COMPLEJIDAD DINÁMICA REGIONAL ---

La vulnerabilidad frente a los ciclones extratropicales ha quedado recientemente expuesta en el sur de Brasil, donde \citet{bartolomei2024extremos} documentan la letalidad causada por estos sistemas invernales. Sin embargo, la caracterización precisa de su estructura dinámica en el Hemisferio Sur ha estado históricamente limitada por la resolución espacial de los datos disponibles. Esta limitación no es menor: \citet{priestley2022improved} demuestran que los modelos de clima subestiman sistemáticamente los vientos extremos asociados a ciclones al no resolver explícitamente sus estructuras de mesoescala, lo que hace indispensable el uso de productos de alta resolución como ERA5 para capturar adecuadamente la organización espacial de los campos de superficie. En el Hemisferio Norte, investigaciones recientes de \citet{gentile2025response} anticipan, mediante modelaciones de alta resolución que el sector cálido de los ciclones extratropicales más intensos se intensificará, incrementando significativamente los vientos y la precipitación extrema durante sus etapas de desarrollo.

A pesar de la relevancia dinámica de estos sistemas, la literatura sobre el Atlántico Sur ha priorizado históricamente el análisis de la climatología de trayectorias y la frecuencia de ciclogénesis. Si bien referentes como \citet{gramcianinov2020analysis} e investigaciones recientes como \citet{padilhareinke2026characterization} han caracterizado con robustez la distribución espacio-temporal de estos sistemas, persiste un vacío: la comprensión de cómo se estructura espacialmente los campos de viento y precipitación durante la evolución del ciclon, y cómo esta organización interna transforma su configuración a través de las distintas fases del ciclo de vida. Esta perspectiva estructural es necesaria en sudamerica, en la región trabajos de \citet{gozzo2013air} y \citet{reboita2022from} han revelado que los ciclones frecuentemente exhiben evoluciones atípicas que no pueden ser capturadas por métricas simples de intensidad central. Esta heterogeneidad dinámica se traduce en una distribución de severidad altamente asimétrica; investigaciones de \citet{bitencourt2010relating} y \citet{cardoso2022synoptic} evidencian que los máximos de viento en superficie no se distribuyen uniformemente alrededor del vórtice, sino que dependen críticamente de la latitud de génesis y la madurez del sistema.

En esa misma linea, \citet{eisenstein2023identification}, aunque para el hemisferio norte, advierten que los vientos más intensos suelen estar asociados a estructuras de mesoescala muy localizadas, cuya posición varía durante el ciclo de vida. Así, autores como \citet{zscheischler2018future} señalan que existe un interes en la ocurrencia temporal y local de los máximos de eventos compuestos de viento y precipitación, la ocurrencia simultanea de estas variables superficiales como efectos combinados. 

En el Atlántico Sur, los ciclones con estructuras híbridas están frecuentemente vinculados a condiciones de tiempo adverso \citep{marrafon2022classificacao}. Sin embargo, la evolución conjunta de estas variables a través de las fases de vida del ciclón permanece escasamente documentada. Recientemente \citet{han2025system} enfatiza que se requieren marcos de clasificación que vinculen explícitamente la evolución del sistema con sus manifestaciones extremas de viento y precipitación, a esto nos referimos como los eventos compuestos mencionados anteriormente, un desafío abordable en Sudamérica ahora que se cuenta con metodologías objetivas como la de \citet{coutodesouza2024new} para la clasificacion de las fases de vida del ciclon.

% --- PÁRRAFO 3: EL FORZANTE OROGRÁFICO Y LA REGIONALIZACIÓN (SBR, LPB, ARG) ---

Abordar este desafío en el Atlántico Sur exige, sin embargo, comprender los mecanismos regionales que modulan la génesis y la estructura interna de estos sistemas. En este dominio, el continente sudamericano introduce una perturbación orográfica fundamental que condiciona tanto la localización de la ciclogénesis como la organización espacial de los campos de superficie. La Cordillera de los Andes interactúa mecánicamente con el flujo de los oestes, forzando la compresión de la columna de vorticidad a barlovento y su estiramiento vertical a sotavento, un mecanismo físico descrito por \citet{gan1994influence} que favorece la formación de centros de baja presión al este de la cordillera, proceso conocido como ciclogénesis de sotavento.

Este forzante topográfico, la fuerte baroclinicidad costera y los gradientes de temperatura superficial del mar, consolida tres focos principales de actividad ciclogenética en la costa este de Sudamérica \citep{reboita2026meteorology}. 

\label{sec:definicion_regiones}%
Adoptando la regionalización propuesta por \citet{gramcianinov2019properties}, este estudio se centra en tres regiones clave: la región Sur-Sudeste de Brasil (\textit{South Brazil}, SBR), la cuenca de descarga del Río de la Plata (\textit{La Plata Basin}, LPB, por sus siglas en inglés), y la costa sureste de Argentina (\textit{Argentina}, ARG). Estas regiones poseen una correspondencia espacial directa con las regiones de genesis (RG) clásicamente definidos por \citet{reboita2010south} y revisados por \citet{reboita2026meteorology}: (1) la región SBR (análoga a RG1), favorecida por el efecto de sotavento de los Andes y el chorro de bajos niveles (SALLJ), donde \citet{reboita2022from} documentan frecuentes sistemas híbridos y transiciones tropicales; (2) la región LPB (correspondiente a RG2), caracterizada por ciclogénesis explosivas e influenciada por el transporte de humedad del SALLJ, resultando determinante para la inestabilización de la atmósfera \citep{crespo2020potential}; y (3) la región ARG (equivalente a RG3), dominada por la baroclinia del frente polar, donde la dinámica de chorro polar y la propagación de ondas de Rossby constituyen los principales mecanismos de desarrollo \citep{simmonds2000variability}.

Esta estratificación geográfica en el diseño se debe a que estos focos responden a regímenes dinámicos distintos. Evidencia de ello son las proyecciones de \citet{dejesus2021multimodel}, quienes demuestran que estas regiones presentan sensibilidades opuestas ante el forzamiento climático: mientras que la región austral (ARG) experimentará un aumento en la frecuencia de ciclones debido al desplazamiento de los oestes hacia el polo, se proyecta una disminución en los focos subtropicales (SBR y LPB), sugiriendo que son sistemas gobernados por mecanismos independientes. Proyecciones con modelos regionales de alta resolución confirman que, aunque se espera una reducción en la frecuencia ciclónica en SBR y LPB hacia fines de siglo, los sistemas que persistan exhibirán una intensificación de los flujos de humedad y la advección de calor en capas bajas \citep{gramcianinov2024early}, lo que intensificaría la severidad de los eventos en estas regiones.

% --- PÁRRAFO 4: INGREDIENTES DINÁMICOS Y EL ROL DE LA HUMEDAD ---

La variabilidad de la actividad ciclónica subtropical es modulada por la corriente en chorro y su interacción con las capas bajas. Al cruzar los Andes, segun \citet{vera2002cold}, y realinearse en el lado de sotavento, generan convergencia de vientos en niveles bajos y divergencia en altura, lo que intensifica los ciclones y provoca precipitaciones intensas en el sureste de Sudamérica. Simultáneamente, el transporte meridional de humedad desde los trópicos, mediado por el Chorro de Bajos Niveles de Sudamérica (SALLJ, por sus siglas en inglés), juega un rol determinante en la alimentación de la precipitación extrema. Como señalan \citet{reboita2022from}, la advección de aire cálido y húmedo en el flanco oriental del ciclón no solo inestabiliza la atmósfera, sino que actúa como una fuente diabática que, al liberar calor latente, refuerza la intensidad del sistema. Sin embargo, este acoplamiento termodinámico introduce una complejidad crítica, \citet{sinclair2023relationship} advierten que la relación entre la intensidad del sistema (vorticidad) y su respuesta (precipitación) no es lineal, ya que un ciclón puede ser dinámicamente intenso pero generar poca lluvia si carece de este suministro de humedad. Este hecho físico subraya que los máximos de severidad no necesariamente coinciden con el mínimo de presión, reforzando la necesidad de caracterizar la huella espacial de viento y precipitación.

% --- PÁRRAFO 5: DESACOPLAMIENTO ESPACIAL Y MAGNITUD DE LOS EXTREMOS ---

Un desafío en la caracterización de estos sistemas es su gran complejidad estructural y asimetría espacial; esta variabilidad en la forma y tamaño de los ciclones extratropicales genera incertidumbres en la determinación de su intensidad en superficie, especialmente en regiones con topografía compleja \citep{neu2013intercomparison}. Mientras que la literatura clásica vinculaba la fuerza de los ciclones exclusivamente con su intensidad central de presion, actuales investigaciones demuestran un desacoplamiento sistemático entre la profundidad del centro y la magnitud de sus extremos asociados. \citet{bitencourt2010relating} y \citet{cardoso2022synoptic} evidencian que los máximos de viento no se distribuyen uniformemente alrededor del centro, sino que se concentran en sectores específicos (como el cuadrante noreste en la región subtropical). Esta desconexión sugiere que la métrica de presión es insuficiente para capturar la severidad del evento, por otro lado, \citet{cardoso2022synoptic} documentan que los ciclones intensos según presión mínima (percentil 10 de la Presión media a nivel del mar) presentan trayectorias sistemáticamente más alejadas de la costa que los intensos según viento máximo, evidenciando que esta medida no garantiza la zona de impacto ni severidad en superficie.

En la literatura clásica del Hemisferio Norte, el flujo ascendente del Cinturón Transportador Cálido (WCB, por sus siglas en ingles) % Warm Conveyor Belt
se describe típicamente hacia el noreste \citep{blanchard2021warm}. El WCB no solo produce precipitación a gran escala, sino que frecuentemente alberga núcleos de convección intensa embebida \citep{oertel2021warm}, contribuyendo a la producción de precipitación extrema desacoplada del centro del ciclón \citep{naud2020evaluation}. Como muestran \citet{chen2025characteristics}, la caracterización física de estos eventos emerge de la co-ocurrencia temporal y local de extremos de viento y precipitación, una superposición que no es capturada al analizar cada variable de forma individual. De esta forma, se hace necesario una caracterización de la huella espacial, que permita cuantificar la distribución de estas variables en superficie y de sus máximos. 

% --- PÁRRAFO 6  Pregunta de investigación e Hipótesis ---

Ante esta complejidad estructural y temporal, se plantea la interrogante de esta investigación: ¿cómo evoluciona espacialmente los campos de viento y precipitación en los ciclones extratropicales, y qué configuraciones caracterizan esta organización durante las fases incipiente, intensificación, madurez y decaimiento? 
Se postula que la estructura de viento y precipitación presenta patrones de asimetría sistemática que varían entre fases evolutivas.

% --- PÁRRAFO 7: EVOLUCIÓN TEMPORAL Y CLASIFICACIÓN DINÁMICA (FASES) ---

Esta complejidad estructural no es estática; evoluciona drásticamente a lo largo del ciclo de vida del sistema. Frente al modelo clásico, que presenta una oclusión simétrica, los ciclones del Atlántico Sur suelen seguir el modelo de Fractura Frontal propuesto por \citet{shapiro1990life}, el cual describe la evolución de sistemas que experimentan intensificación rápida con desarrollo de seclusión cálida \citep{shapiro1999bridge}. En este marco, el WCB y su respectiva corriente en chorro de bajo nivel se organizan en torno al vórtice \citep{browning1986conceptual}, configurando la estructura tridimensional del sistema precipitante y determinando la distribución espacial de los máximos de viento y precipitación.

Históricamente, la clasificación de estos eventos se basó en el diagrama de espacio de fase de \citet{hart2003cyclone}, el cual discrimina sistemas según su asimetría térmica y su núcleo (frío/caliente). Sin embargo, este enfoque presenta limitaciones operativas en el Atlántico Sur, donde \citet{reboita2022from} han documentado frecuentes transiciones rápidas y sistemas híbridos que desafían las categorías binarias clásicas. Para caracterizar la evolución física de la estructura interna, esta investigación adopta la segmentación objetiva desarrollada por \citet{coutodesouza2024new}, quienes mediante el algoritmo \textit{CycloPhaser} han definido cuatro fases dinámicas discretas —incipiente, intensificación, madurez y decaimiento— \citep{deSouza2025CycloPhaser}, estableciendo una base física para analizar la evolución de los campos de superficie a través del ciclo de vida del sistema. El uso de esta base de datos clasificada permite, por primera vez, investigar cómo se transforma la distribución espacial de las variables durante sus fases del ciclón, una estrategia alineada con otras nuevas propuestas globales como \textit{SyCLoPS} de \citet{han2025system} en el Atlantico Norte, quienes adicionalmente enfatizan la urgencia de vincular la evolución física del sistema con sus patrones de severidad asociados.

\label{subsec:intro_estrategia_estadistica}%
Para caracterizar esta complejidad estructural, se requieren técnicas estadísticas que trasciendan el análisis univariado clásico y capturen la covariabilidad entre campos físicos distintos. Si bien enfoques recientes como el de \citet{eisenstein2023identification} han demostrado que las configuraciones extremas del viento pueden clasificarse en tipos discretos mediante aprendizaje automático, la variabilidad espacial continua de la estructura ciclónica ---particularmente la coherencia entre viento y precipitación--- requiere herramientas de reducción dimensional que preserven la interpretabilidad física. En este sentido, el análisis de Funciones Ortogonales Empíricas Multivariadas (MEOF) se establece como el marco apropiado \citep{hannachi2007empirical}: A diferencia del análisis univariado de EOF aplicado separadamente a cada campo, la técnica MEOF permite extraer modos de variabilidad inherentemente acoplados entre variables físicamente distintas \citep{liang2018multivariate}, revelando la organización espacial coherente donde la circulación del viento y los núcleos de precipitación covarían durante la evolución del sistema.

La aplicación de este marco estadístico sobre las fases del ciclon previamente clasificadas permitirá construir "prototipos estructurales", sintetizando la estructura típica del ciclón para cada región de génesis y fase evolutivo, proponiendo un analisis alternativo a la visión estática de los promedios climatológicos.

% --- PÁRRAFO 8: DATOS Y HERRAMIENTAS (ERA5 Y TRACKING POR VORTICIDAD) ---

La viabilidad de esta caracterización estructural reside en la disponibilidad del reanálisis ERA5 (European Centre for Medium-Range Weather Forecasts Reanalysis Fifth Generation) \citep{hersbach2020era5}, cuya resolución espacial ($\sim$31 km) y temporal (horaria) permite resolver los gradientes de presión y los flujos de humedad que generaciones anteriores de datos suavizaban. Validaciones específicas para el Atlántico Sur realizadas por \citet{gramcianinov2020analysis} confirman que ERA5 ofrece una representación superior de la intensidad de los vientos superficiales y la estructura de los ciclones en comparación con productos previos.
No obstante, variables de ERA5 en superficie asociadas a ciclones demuestran que, si bien el producto representa el comportamiento promedio de los sistemas, tiende a subestimar la magnitud de los extremos más intensos \citep{chen2024evaluation}.

Para capitalizar esta precisión, este estudio utiliza la base de datos de trayectorias de \citet{coutodesouza2024new}, la cual se fundamenta en el criterio de la vorticidad relativa en 850 hPa ($\zeta_{850}$). Esta elección metodológica posee un consenso transversal en la literatura: para resolver el problema de aislar la circulación del sistema, la recomendación de \citet{sinclair1997objective} es utilizar $\zeta_{850}$ como variable de seguimiento, una práctica establecida clásicamente por \citet{sinclair1994objective} y tambien usada por \citet{padilhareinke2024objective} como el estándar más robusto para el Hemisferio Sur. La literatura reciente utiliza $\zeta_{850}$ para identificar con precisión las regiones de ciclogénesis \citep{gramcianinov2019properties}, mientras otros estudios de validación demuestran su superioridad sobre campos de presión \citep{padilhareinke2024objective}, siendo incluso utilizado en estudios recientes del Hemisferio Norte \citep{chen2025characteristics}. A diferencia de los algoritmos basados exclusivamente en la presión mínima, el tracking por vorticidad permite el análisis del ciclo de vida de los eventos, capturando incluso aquellos sistemas rápidos que los algoritmos isobáricos tienden a omitir. Esta combinación de datos de alta fidelidad con una detección dinámica sensible constituye la base técnica necesaria para alimentar el análisis multivariado (MEOF) propuesto, permitiendo cuantificar finalmente la evolución de la huella espacial de los eventos compuestos en la región.

\section{Objetivos del Proyecto de Investigación}

\subsection{Objetivo General}
Caracterizar la evolución de la estructura espacial del viento a 10 metros y precipitación durante las fases ciclo de vida de ciclones extratropicales en el Atlántico Sur, caracterizando sus patrones de variabilidad y estimación de densidad probabilística de máximos, con estratificación por región de génesis.

\subsection{Objetivos Específicos}
\begin{enumerate}
    \item Consolidar una base de datos lagrangiana que asocie los campos horarios de alta resolución de viento y precipitación (ERA5) con las trayectorias y fases dinámicas (incipiente, intensificación, madurez y decaimiento) de ciclones extratropicales del Atlántico Sur, estratificando la muestra por región de génesis (SBR, LPB y ARG) y por intensidad dinámica (Población Global y subconjunto p90).
    \item Cuantificar la distribución estadística y la asimetría espacial de los máximos de viento a 10 m y precipitación mediante estimación de densidad por núcleos, caracterizando tanto la distribución de probabilidad de sus magnitudes (PDF unidimensional) como la localización preferencial de dichos máximos respecto al centro del ciclón (PDF bidimensional).
    \item Identificar los modos dominantes de variabilidad espacial de los campos de viento y precipitación mediante Funciones Ortogonales Empíricas univariadas (EOF) y multivariada (MEOF), evaluando el patron estructural individual de cada variable como la covariabilidad acoplada.
    \item Construir prototipos estructurales sintéticos por fase evolutiva y región de génesis, integrando los campos medios condicionados al modo dominante de covariabilidad (MEOF) con las densidades probabilísticas de localización de máximos (PDFe), para cuantificar la magnitud característica de viento y precipitación.    
\end{enumerate}